\chapter{АШИГЛАЛТЫН ХУВИЛБАРУУДЫН ДЭЛГЭРЭНГҮЙ ЖАГСААЛТ}

Энэ хавсралтад төслийн хүрээнд тодорхойлогдсон хэрэглэгчийн ашиглалтын хувилбаруудыг (use case) дэлгэрэнгүй тайлбарлана. Эдгээр нь системийн хөгжүүлэлт, туршилтын сценари боловсруулах үндэс болно. Хувилбар бүрийг тодорхойлохдоо хэрэглэгчийн хүсэлтийн жишээ, системийн хүлээгдэж буй хариу үйлдэл, шаардлагатай MCP хэрэгслүүд зэргийг тусгасан.

\section{Оюутны талын үндсэн хувилбарууд}

Эдгээр нь одоогийн MCP хэрэгслүүдээр дэмжигдэх боломжтой read-only хувилбарууд юм.

\subsection{Хуваарь харах}

\textbf{Хэрэглэгчийн хүсэлт:} ``Миний энэ долоо хоногийн хичээлийн хуваарийг харуулаач.''

\textbf{Системийн хариу үйлдэл:} Систем нь \texttt{get\_schedule} хэрэгслийг дуудаж, хэрэглэгчийн хичээлийн хуваарийг авна. Дараа нь өгөгдлийг боловсруулж, долоо хоногийн өдөр бүрээр ангилсан хуваарийг хэрэглэгчид ойлгомжтой хэлбэрээр харуулна. Хэрэв шалгалтын хуваарь байвал түүнийг тусад нь онцолж мэдээлнэ.

\textbf{Ашиглах MCP хэрэгсэл:} \texttt{get\_schedule}

\subsection{Дүн харах}

\textbf{Хэрэглэгчийн хүсэлт:} ``Миний өнгөрсөн семестрийн дүнг харуулаач. Голч хэд вэ?''

\textbf{Системийн хариу үйлдэл:} Систем нь \texttt{get\_academic\_record} хэрэгслийг дуудаж, хэрэглэгчийн академик бүртгэлийг авна. Өнгөрсөн семестрийн хичээлүүд, тэдгээрийн дүн, нийт голч (GPA) зэргийг тооцоолж харуулна. Хэрэв хэрэглэгч тодорхой хичээлийн талаар асуувал тухайн хичээлийн дүнг онцлон харуулна.

\textbf{Ашиглах MCP хэрэгсэл:} \texttt{get\_academic\_record}

\subsection{Төлбөрийн мэдээлэл}

\textbf{Хэрэглэгчийн хүсэлт:} ``Миний төлбөрийн үлдэгдэл хэд вэ? Хэзээ төлөх ёстой вэ?''

\textbf{Системийн хариу үйлдэл:} Систем нь \texttt{get\_financial\_overview} хэрэгслийг дуудаж, хэрэглэгчийн санхүүгийн мэдээллийг авна. Одоогийн үлдэгдэл, төлбөрийн хугацаа, өмнөх төлөлтүүдийн түүх зэргийг харуулна. Хэрэв төлбөрийн хугацаа ойрхон байвал анхааруулга өгнө.

\textbf{Ашиглах MCP хэрэгсэл:} \texttt{get\_financial\_overview}

\subsection{Мэдэгдэл харах}

\textbf{Хэрэглэгчийн хүсэлт:} ``Ямар нэгэн шинэ мэдэгдэл байна уу?''

\textbf{Системийн хариу үйлдэл:} Систем нь \texttt{get\_notifications} хэрэгслийг дуудаж, хэрэглэгчийн мэдэгдлүүдийг авна. Уншаагүй мэдэгдлүүдийг эхэлж харуулна. Мэдэгдлийн төрөл (захиргааны, хичээлийн, санхүүгийн гэх мэт)-өөр ангилж болно.

\textbf{Ашиглах MCP хэрэгсэл:} \texttt{get\_notifications}

\subsection{Ерөнхий мэдээлэл}

\textbf{Хэрэглэгчийн хүсэлт:} ``Миний оюутны мэдээллийг харуулаач.''

\textbf{Системийн хариу үйлдэл:} Систем нь \texttt{get\_student\_overview} хэрэгслийг дуудаж, хэрэглэгчийн ерөнхий мэдээллийг авна. Нэр, оюутны дугаар, мэргэжил, курс, статус зэргийг харуулна.

\textbf{Ашиглах MCP хэрэгсэл:} \texttt{get\_student\_overview}

\subsection{Олон мэдээлэл нэг дор}

\textbf{Хэрэглэгчийн хүсэлт:} ``Маргааш ямар хичээл байгаа, мөн дүнгийн мэдээлэл хэлээч.''

\textbf{Системийн хариу үйлдэл:} Систем нь хоёр хэрэгслийг дуудна: \texttt{get\_schedule} болон \texttt{get\_academic\_record}. Хоёр эх үүсвэрийн мэдээллийг нэгтгэж, маргаашийн хуваарь болон дүнгийн мэдээллийг хамтад нь харуулна.

\textbf{Ашиглах MCP хэрэгсэл:} \texttt{get\_schedule}, \texttt{get\_academic\_record}

\section{Өгөгдөл экспортлох хувилбарууд}

Эдгээр нь одоогийн MCP хэрэгслүүдийн гаралтыг файл хэлбэрт хөрвүүлэх хувилбарууд юм.

\subsection{Дүн CSV экспорт}

\textbf{Хэрэглэгчийн хүсэлт:} ``Миний бүх дүнг CSV файл болгож өгөөч.''

\textbf{Системийн хариу үйлдэл:} Систем нь \texttt{get\_academic\_record} хэрэгслийг дуудаж, бүх академик бүртгэлийг авна. Дараа нь өгөгдлийг CSV формат руу хөрвүүлж, хэрэглэгчид татаж авах боломж олгоно. CSV файлд хичээлийн нэр, кредит, дүн, семестр зэрэг багтана.

\textbf{Ашиглах MCP хэрэгсэл:} \texttt{get\_academic\_record}

\textbf{Тэмдэглэл:} Энэ нь read-only өгөгдөл дээр суурилсан хувиргалт бөгөөд mutation биш. AI загвар нь өгөгдлийг форматлах чадвартай.

\subsection{Хуваарь календарь руу}

\textbf{Хэрэглэгчийн хүсэлт:} ``Миний хичээлийн хуваарийг Google Calendar-т нэмж өгөөч.''

\textbf{Системийн хариу үйлдэл:} Систем нь \texttt{get\_schedule} хэрэгслийг дуудаж, хуваарийг авна. Дараа нь ерөнхий MCP өргөтгөлөөр (жишээ нь Google Calendar MCP) календарт event үүсгэнэ.

\textbf{Ашиглах MCP хэрэгсэл:} \texttt{get\_schedule}, Google Calendar MCP (төлөвлөсөн)

\textbf{Тэмдэглэл:} Энэ хувилбар нь ерөнхий MCP өргөтгөл шаарддаг бөгөөд дараагийн шатанд хэрэгжинэ.

\section{Ерөнхий MCP өргөтгөлийн хувилбарууд}

Эдгээр нь SISI-ийн өгөгдлийг гадаад системүүдтэй холбох хувилбарууд юм. Ирээдүйн хэрэгжилтэд багтана.

\subsection{Даалгавар reminder}

\textbf{Хэрэглэгчийн хүсэлт:} ``Миний даалгавруудын deadline-ийг reminder болгож тохируулаач.''

\textbf{Системийн хариу үйлдэл:} Систем нь SISI-ээс даалгаврын мэдээлэл авч, reminder MCP (жишээ нь Todoist, Notion) ашиглан хэрэглэгчийн сонгосон апп-д reminder үүсгэнэ.

\textbf{Ашиглах MCP хэрэгсэл:} SISI даалгавар хэрэгсэл, Reminder MCP (төлөвлөсөн)

\subsection{Тэмдэглэл хадгалах}

\textbf{Хэрэглэгчийн хүсэлт:} ``Энэ мэдээллийг Notion-д хадгалаач.''

\textbf{Системийн хариу үйлдэл:} Систем нь хэрэглэгчийн зааж өгсөн мэдээллийг Notion MCP ашиглан Notion database эсвэл page-д хадгална.

\textbf{Ашиглах MCP хэрэгсэл:} Notion MCP (төлөвлөсөн)

\section{Багшийн талын хувилбарууд}

Эдгээр нь багшийн порталтай холбогдсон хувилбарууд бөгөөд эцсийн дипломын ажилд хэрэгжинэ. Mutation үйлдэл агуулдаг тул аюулгүй байдлын нэмэлт шаардлагатай.

\subsection{Excel-ээс дүн оруулах}

\textbf{Хэрэглэгчийн хүсэлт:} ``Энэ Excel файлаас оюутнуудын дүнг системд оруулаач.''

\textbf{Системийн хариу үйлдэл:} Систем нь дараах алхмуудыг гүйцэтгэнэ. Эхлээд Excel файлыг уншиж, мөр бүрийн өгөгдлийг задлана. Дараа нь мөр бүрийг шалгана: оюутны код зөв эсэх, дүн зөвшөөрөгдсөн хүрээнд байгаа эсэх, хичээлийн код таарч байгаа эсэх. Алдаатай мөрүүдийг жагсааж, багшид харуулна. Багш баталгаажуулсны дараа зөв мөрүүдийг системд оруулна. Эцэст нь үр дүнгийн тайлан гаргана: амжилттай оруулсан, алдаатай, алгассан мөрүүдийн тоо.

\textbf{Ашиглах MCP хэрэгсэл:} Багшийн портал MCP (төлөвлөсөн), Excel parsing

\textbf{Аюулгүй байдал:} Энэ нь mutation үйлдэл тул хэд хэдэн хамгаалалт шаардлагатай. Оруулахын өмнө бүх өгөгдлийг багшид харуулж баталгаажуулалт авах шаардлагатай. Мөр бүрийн validation хийх хэрэгтэй. Audit trail бичих буюу хэн, хэзээ, ямар өгөгдөл оруулсныг бүртгэх шаардлагатай. Rollback буюу буцаах боломжтой байх хэрэгтэй.

\subsection{Дүнгийн тайлан}

\textbf{Хэрэглэгчийн хүсэлт:} ``Миний ангийн дүнгийн тархалтыг харуулаач.''

\textbf{Системийн хариу үйлдэл:} Систем нь багшийн порталаас ангийн дүнгийн өгөгдлийг авч, статистик тооцоолно: дундаж, медиан, дээд/доод дүн, дүнгийн тархалт (histogram). Энэ мэдээллийг багшид ойлгомжтой хэлбэрээр харуулна.

\textbf{Ашиглах MCP хэрэгсэл:} Багшийн портал MCP (төлөвлөсөн)

\section{Нарийвчилсан харилцааны хувилбарууд}

Эдгээр нь хэрэглэгчийн асуултад илүү нарийвчилсан хариу шаардах хувилбарууд юм.

\subsection{Хичээл санал болгох}

\textbf{Хэрэглэгчийн хүсэлт:} ``Дараагийн семестрт ямар хичээл үзэх боломжтой вэ? Урьдчилсан шаардлагыг хангасан хичээлүүдийг харуулаач.''

\textbf{Системийн хариу үйлдэл:} Систем нь \texttt{get\_academic\_record} болон хичээлийн мэдээллийн хэрэгслийг ашиглан хэрэглэгчийн үзсэн хичээлүүд, тэдгээрийн дүнг авна. Дараа нь урьдчилсан шаардлагын (prerequisite) логикоор үзэх боломжтой хичээлүүдийг тооцоолж санал болгоно.

\textbf{Ашиглах MCP хэрэгсэл:} \texttt{get\_academic\_record}, хичээлийн каталог хэрэгсэл (төлөвлөсөн)

\subsection{Төгсөлтийн шаардлага}

\textbf{Хэрэглэгчийн хүсэлт:} ``Төгсөхөд хэдэн кредит дутуу байна вэ? Ямар хичээл заавал үзэх ёстой вэ?''

\textbf{Системийн хариу үйлдэл:} Систем нь хэрэглэгчийн академик бүртгэл болон мэргэжлийн төгсөлтийн шаардлагыг харьцуулж, дутуу кредит, заавал үзэх хичээлүүдийг тооцоолно.

\textbf{Ашиглах MCP хэрэгсэл:} \texttt{get\_academic\_record}, \texttt{get\_student\_overview}, төгсөлтийн шаардлагын хэрэгсэл (төлөвлөсөн)

\section{Тэмдэглэл}

Дээрх use case-уудаас зарим нь одоогийн хэрэгжилтээр дэмжигдэж байгаа бол зарим нь дараагийн шатны ажлуудтай холбоотой.

Одоо боломжтой хувилбарууд нь хуваарь харах, дүн харах, төлбөрийн мэдээлэл, мэдэгдэл харах, ерөнхий мэдээлэл, олон мэдээлэл нэг дор, CSV экспорт (AI форматлалт) зэрэг юм.

Дараагийн шатанд хэрэгжих хувилбарууд нь ерөнхий MCP өргөтгөлүүд (календарь, Notion), хичээл санал болгох, төгсөлтийн шаардлага зэрэг юм.

Эцсийн дипломд хэрэгжих хувилбарууд нь багшийн порталын бүх хувилбарууд, Excel дүн оруулах, mutation баталгаажуулалт зэрэг юм.

Туршилтын шатанд эдгээр use case-уудыг ашиглан системийн амжилтын хувь, хариу өгөх хугацаа, хэрэглэгчийн сэтгэл ханамж зэргийг хэмжинэ.
